\documentclass{abstract}

% Bibliography
\usepackage[
    backend=biber,
    style=authoryear,
    doi=true,
    sorting=none,
    sortcites=true,
    maxbibnames=2,
    minbibnames=1,
    maxcitenames=2,
    mincitenames=1,
    citestyle=numeric-comp,
    giveninits=true,
    isbn=false,
    date=year
]{biblatex}
\addbibresource{abstract.bib}

% For email linking
\usepackage{hyperref}

\title{Effect of heat flux on combustion of different wood species}

\author{Ian Kim},

\keywords{}
\begin{document}
\maketitle
\section{Abstract}
The main point of this NIST research effort is to look at the cone calorimetry, which is the standard instrument for large-scale fire testing used to measure the flammability of materials by measuring the heat release rate (HRR) from the smoke and oxygen emitted. 
%"look at the cone calorimetry"? rephrase
%% This is not a full-scale test instrument, rather bench-scale. 
% When introducing the cone, more detail should be provided on configuration/how sample's are heated/what this represents... 
And from 1982-1991, NIST conducted over 1000+ cone calorimetry tests looking at the heat release rate of different materials. However, these tests were either handwritten or scattered across different databases. 
%Tests have been run since the 80s, we're looking at 40+ years of data, not 9 years. What materials (in general?)
% Where do we have handwritten data? What other databases? I'm not sure these statements are correct.
%Say where this data actually was published instead...
%%Consider this abstract from the view of someone unfamiliar with fire. We should intro in 2-3 sentence what the cone is, how it works, what data is available (where/how), and what's missing (i.e., the point of this effort, how can we to access this in more useful way)..

Therefore, our mission this spring with the ENFP 429 research class is to scan, review, and centralize this data so that anyone in the world can see how heat release changes from different materials and apply our data to future burn tests and fire protection systems. 
%This is (a) not what you're doing - you're not scanning.. please be technically correct but capture the big picture and (b) the "impact/how this is used.. this isn't quite right either. I think you're mixing what the cone measures, how that data can be used/interpretted, and what new things our DB will enable. Perhaps consider reviewing this simple overview before we catch up next week: https://www.nist.gov/laboratories/tools-instruments/cone-calorimeter

Example citation \cite{babrauskas1982development, brown1988cone,emil1989flammability}


\section{References}
\printbibliography[heading=none]

\end{document}